\chapter{Introdução}
\label{cap:introducao}

O diabetes mellitus é uma doença crônica de prevalência crescente em todo o mundo e representa um importante problema de saúde pública. Segundo a \citeonline[p.~41]{idfAtlas2025}, cerca de 589 milhões de adultos entre 20 e 79 anos viviam com diabetes em 2024, o que corresponde a 11,1\% das pessoas dessa faixa etária, e a estimativa é de que esse número chegue a 853 milhões em 2050. O Brasil ocupa o sexto lugar nessa lista, com cerca de 16,6 milhões de adultos com a doença \cite[p.~47]{idfAtlas2025}. Ao longo do tempo, o diabetes pode danificar os vasos sanguíneos do coração, dos olhos, dos rins e dos nervos, e constitui uma causa importante de cegueira, insuficiência renal, infarto e amputações \cite{who2026diabetes}.

No olho, o diabetes pode causar a retinopatia diabética (RD), uma complicação microvascular comum e específica da doença, que afeta os pequenos vasos da retina \cite{brasil2021pcdtrd,malerbi2026rd}. Em uma revisão com dados de diversos países, \citeonline{teo2021global} estimaram que cerca de 22\% das pessoas com diabetes apresentam algum grau de RD, o que correspondia a 103 milhões de pessoas em 2020, com previsão de 160 milhões em 2045. No Brasil, segundo o Protocolo Clínico e Diretrizes Terapêuticas (PCDT) do Ministério da Saúde, a RD atinge entre 24\% e 39\% das pessoas com diabetes, e figura entre as principais causas de perda visual na faixa etária de 20 a 75 anos \cite{brasil2021pcdtrd}. Além de frequente, a RD pode evoluir de forma silenciosa, uma vez que, no início da doença, a visão do paciente costuma permanecer preservada \cite{brasil2021pcdtrd}. Segundo a diretriz da Sociedade Brasileira de Diabetes (SBD), quando a RD é diagnosticada tardiamente e alcança as formas graves, ela passa a ser a principal causa de perda de visão evitável entre adultos em idade produtiva \cite{malerbi2026rd}. Em 1999, o consenso da SBD e do Conselho Brasileiro de Oftalmologia (CBO) já apontava que a maior parte dos casos de cegueira ligados à RD pode ser evitada com medidas adequadas \cite{gross1999consenso}, o que torna fundamental o rastreamento das pessoas com diabetes.

O rastreamento é realizado por meio do exame do fundo do olho. A SBD orienta que a primeira avaliação oftalmológica seja feita após a dilatação da pupila e inclua, no mínimo, a retinografia ou o mapeamento de retina, além da biomicroscopia de fundo \cite{malerbi2026rd}, e o PCDT cita a fotografia da retina como uma forma de empregar a telessaúde em locais com pouco ou nenhum acesso a oftalmologista \cite{brasil2021pcdtrd}. Entretanto, esse acesso ainda é desigual no país. Segundo o \citeonline[p.~7]{cbo2021censo}, os postos de trabalho de oftalmologistas em atividade estão em apenas 1.689 das 5.570 cidades brasileiras, e cerca de 43,3 milhões de habitantes vivem em 3.880 municípios sem esse profissional, de modo que grande parte dos municípios brasileiros não conta com esse especialista. Um levantamento mais recente, feito com outra fonte de dados, o Cadastro Nacional de Estabelecimentos de Saúde (CNES) de 2024, encontrou pelo menos um oftalmologista em apenas 23,5\% dos municípios, com desigualdade ainda maior no setor público \cite{gomes2026oftalmologistas}. A falta de profissionais treinados também aparece no relatório mundial sobre visão da Organização Mundial da Saúde (OMS) como um dos maiores desafios para ampliar o cuidado com os olhos \cite[p.~36]{who2019vision}. Além da falta de acesso, a leitura das imagens é subjetiva, e o grau dado à mesma imagem pode variar entre avaliadores, inclusive entre especialistas \cite{krause2018grader}.

Por meio da retinografia, é possível observar lesões que ajudam a avaliar a RD, tais como microaneurismas, hemorragias, exsudatos duros e exsudatos moles \cite{porwal2018idrid}. Os exsudatos duros correspondem a acúmulos de gordura e proteína que vazam de vasos lesados e se apresentam como pontos amarelados e brilhantes, enquanto os exsudatos moles, também chamados de manchas algodonosas, são manchas esbranquiçadas associadas à isquemia da retina \cite{negi2003eye}. Os exsudatos duros têm um papel especial, pois a sua presença é usada para determinar o risco de edema macular \cite{porwal2018idrid}, e o PCDT inclui exsudatos duros perto do centro da mácula na definição de edema macular clinicamente significativo \cite{brasil2021pcdtrd}. Em programas de rastreamento de larga escala, o edema pode ser avaliado pela detecção desses exsudatos nas imagens \cite{giancardo2012exudate}. Por isso, identificar essas lesões na imagem é importante para a gradação da RD e também para decidir se o paciente deve ser encaminhado ao oftalmologista.

Nesse contexto, os sistemas de diagnóstico auxiliado por computador (\textit{Computer-Aided Diagnosis}, CAD) podem ser uma forma sustentável de fazer o rastreamento em larga escala quando faltam profissionais \cite{porwal2020idrid}. O aprendizado profundo, em especial as redes convolucionais, passou a ser uma das principais formas de analisar imagens médicas \cite{litjens2017survey}, e, na RD, já existem resultados expressivos. \citeonline{gulshan2016} treinaram uma rede convolucional que detectou a RD referenciável com área sob a curva ROC de 0,991 no EyePACS-1 e de 0,990 no Messidor-2, e \citeonline{ting2017} validaram um sistema semelhante com 494.661 imagens de populações de várias etnias. O sistema avaliado por \citeonline{abramoff2018pivotal} em clínicas de atenção primária, com sensibilidade de 87,2\% e especificidade de 90,7\%, foi o primeiro sistema autônomo de inteligência artificial autorizado pela agência norte-americana FDA em qualquer área da medicina. No Brasil, \citeonline{malerbi2022smartphone} testaram em Itabuna, na Bahia, uma câmera portátil ligada a um celular junto com um algoritmo de aprendizado profundo, e concluíram que esse tipo de ferramenta pode aumentar a cobertura do rastreamento. Mesmo assim, a OMS ressalta que essas tecnologias ainda precisam de mais estudos em condições reais de uso \cite[p.~75]{who2019vision}.

Nesses estudos, a resposta do sistema é uma decisão sobre a imagem inteira, como o grau da doença ou a necessidade de encaminhar o paciente. A segmentação dos exsudatos acrescenta outra informação, porque mostra onde estão as lesões e permite conferir o resultado na própria imagem. Ela também permite aplicar regras clínicas simples, como a de marcar suspeita de edema macular quando há exsudatos duros a até um diâmetro de disco do centro da fóvea \cite{krause2018grader}. Ao mesmo tempo, existem hoje diversas arquiteturas que podem ser usadas nessas tarefas, como redes convolucionais tradicionais, redes leves pensadas para pouco poder computacional e redes baseadas em \textit{Transformers}. Entre os onze trabalhos analisados no Capítulo~\ref{cap:trabalhos-relacionados}, selecionados por uma busca direcionada e não exaustiva, os que segmentam exsudatos ou classificam a RD costumam propor ou avaliar uma única arquitetura, poucos testam o modelo em uma base diferente da usada no treino e nenhum compara arquiteturas diferentes com teste estatístico. Também não há, entre eles, um estudo que compare redes tradicionais, leves e baseadas em \textit{Transformers} sob o mesmo protocolo. Isso dificulta a escolha de um modelo para uso prático, principalmente em regiões com poucos recursos, onde o custo computacional de cada rede é tão importante quanto o seu desempenho.

Diante disso, este trabalho consiste na realização de um estudo comparativo entre seis arquiteturas de redes de aprendizado profundo, representativas de três famílias, utilizando modelos tradicionais, como VGG16 e ResNet50, modelos leves, como MobileNetV3 e EfficientNet-B0, e modelos baseados em \textit{Transformers}, como ViT e Swin Transformer. As seis redes serão usadas como codificadores de uma mesma U-Net para segmentar os exsudatos duros e moles, com treino no IDRiD e teste também no DDR, e serão treinadas para classificar o grau da RD na escala internacional, de 0 a 4 \cite{wilkinson2003icdr}, com as bases APTOS e EyePACS, e testadas também no Messidor-2. Todas serão treinadas sob o mesmo protocolo, avaliadas também em bases que não foram usadas no treino e comparadas por testes estatísticos e pelo custo computacional. Os modelos mais adequados de cada tarefa serão reunidos em uma solução CAD que mostre as lesões, o grau previsto e um alerta de edema macular, de forma a auxiliar no diagnóstico da retinopatia diabética em regiões com recursos médicos limitados.

\section{Objetivos}
\label{sec:objetivos}

O objetivo geral deste trabalho é comparar seis arquiteturas de aprendizado profundo, representativas das famílias tradicionais, leves e baseadas em \textit{Transformers}, na segmentação de exsudatos duros e moles e na classificação do grau da retinopatia diabética em retinografias, e usar os modelos mais adequados, considerando o desempenho, a generalização e o custo computacional, em uma solução CAD de apoio ao diagnóstico. Os objetivos específicos são:

\begin{alineas}
	\item treinar uma U-Net para segmentar exsudatos duros e moles com seis codificadores diferentes, a VGG16, a ResNet50, a MobileNetV3-Large, a EfficientNet-B0, a ViT-B/16 e a Swin-T, usando a mesma divisão dos dados e a mesma busca de hiperparâmetros para todos;
	\item treinar os mesmos seis modelos para classificar o grau da RD de 0 a 4 e para indicar a RD referenciável;
	\item avaliar os modelos por validação cruzada e também em bases externas, sem nenhum ajuste, para medir a perda de desempenho fora da base de treino;
	\item comparar as arquiteturas por testes estatísticos com correção para várias comparações e pelo custo computacional de cada uma, medido em parâmetros, MACs, tempo de inferência e memória;
	\item desenvolver uma solução CAD com os modelos escolhidos, que mostre as máscaras dos exsudatos, o grau previsto, o alerta de RD referenciável e um alerta de edema macular baseado na posição dos exsudatos duros em relação à fóvea.
\end{alineas}

\section{Organização do trabalho}
\label{sec:organizacao}

Além desta introdução, o trabalho tem mais quatro capítulos. O Capítulo~\ref{cap:fundamentacao-teorica} traz a fundamentação teórica, com a retinopatia diabética e as suas lesões, as arquiteturas avaliadas, a segmentação, as métricas e a comparação estatística entre modelos. O Capítulo~\ref{cap:trabalhos-relacionados} apresenta os trabalhos relacionados e a comparação entre eles e esta proposta. O Capítulo~\ref{chap:metodologia} descreve como o trabalho será feito, com as bases de dados, o pré-processamento, o protocolo de cada tarefa, a análise estatística e a solução CAD. Por fim, o Capítulo~\ref{chap:cronograma} apresenta o cronograma das atividades que serão realizadas no TCC II, com os seus prazos.
